% !PRINTED 420
\weberpage{420}
If $\tau$ is a variable with positive imaginary part, and $v$ likewise a variable, then by § 69 $F_n[v, j(\tau)]$ splits into the linear factors

\[
v - j\left(\frac{c + \partial\,\tau}{a + b\,\tau}\right),
\]

in which $\left(\begin{smallmatrix}a, & b\\ c, & \partial\end{smallmatrix}\right)$ runs through a complete system of representatives of non-equivalent transformations of degree $n$, and consequently $F_n(u, u)$ can be decomposed into the factors

\weq{5}{j(\tau) - j\left(\frac{c + \partial\,\tau}{a + b\,\tau}\right)}

when $u = j(\tau)$ is set.

But from Theorem II., § 114, it follows that, if we set $\tau = \omega$, $k$ of the factors (5) vanish. We can, however, conclude still further that the quotient

\weq{6}{\dfrac{j(\tau) - j\left(\dfrac{c + \partial\,\tau}{a + b\,\tau}\right)}{j(\tau) - j(\omega)}}

remains finite and different from zero for $j(\tau) = j(\omega)$. For if numerator and denominator are differentiated with respect to $\tau$, and $\tau = \omega$ is then set, the limiting value results:

\[
1 - \frac{n}{(a + b\,\omega)^2} = \frac{2\sqrt{D}\,x}{y + \sqrt{D}\,x},
\]

which is different from zero.

The invariant $j(\omega)$ remains unchanged when $\omega$ is replaced by an equivalent number, but it has a different value for every other number $\omega$. It therefore belongs not to the individual form $(A, B, C)$, but to the whole \gesperrt{class} of equivalent forms, and is for that reason called the \gesperrt{class invariant}.

Further, the number $k$ depends not on the coefficients $A, B, C$, but only on $D$ and $n$, and is therefore the same for all primitive forms of discriminant $D$. We now introduce the following function of the variable $u$: let $\omega_1, \omega_2, \ldots, \omega_h$ denote a complete system of non-equivalent numbers of discriminant $D$, that is, the system of the roots of a system of non-equivalent forms $(A, B, C)$ of discriminant $D$. We set

\weq{7}{H_m(u) = [u - j(\omega_1)]\,[u - j(\omega_2)]\,\ldots\,[u - j(\omega_h)],}
%CONT>>
